% AY2012 材料力学 〔I〕（転記）
% 出典: 04_材料力学/original/04_材料力学/2012/2012za.pdf
% 要約: 単純支持・突き出しはり。支点 A, B 間が L, B から右端 C までが L で,
%       右端 C に集中荷重 P が作用する。支持力, SFD/BMD, たわみ基礎式と
%       たわみ角・たわみ曲線・C点たわみ, カスチリアーノの定理によるC点たわみを問う。

\section*{〔I〕}

図に示す単純支持・突き出しはりの右端 C 点に集中荷重 $P$ が負荷される問題に答えよ.
ヤング率（縦弾性係数）を $E$, はりの断面二次モーメントを $I$ とせよ.

\begin{center}
\begin{tikzpicture}[>=Latex]
  % はり本体
  \draw[thick] (0,0) -- (8,0);
  % 座標軸
  \draw[->] (0,0) -- (0,1.4) node[left]{$y$};
  \draw[->] (0,0) -- (1.6,0) node[above]{$x$};
  % 支点 A（ピン）
  \draw (0,0) -- (-0.3,-0.5) -- (0.3,-0.5) -- cycle;
  \draw[pattern=north east lines] (-0.45,-0.7) rectangle (0.45,-0.5);
  \node at (0,-1.0) {A};
  % 支点 B（ローラー）
  \draw (4,0) -- (3.7,-0.5) -- (4.3,-0.5) -- cycle;
  \draw[pattern=north east lines] (3.55,-0.7) rectangle (4.45,-0.5);
  \node at (4,-1.0) {B};
  % 右端 C
  \node at (8,-0.4) {C};
  % 集中荷重 P
  \draw[->,very thick] (8,1.6) -- (8,0.1) node[midway,right]{$P$};
  % 寸法 L, L
  \draw[<->] (0,0.9) -- (4,0.9) node[midway,fill=white]{$L$};
  \draw[<->] (4,0.9) -- (8,0.9) node[midway,fill=white]{$L$};
\end{tikzpicture}
\end{center}

\begin{enumerate}[label=(\arabic*)]
  \item 支持点 A, B に生じる支持力を図示し, つり合い式によりそれらを求めよ.

  \item せん断力 $Q(x)$ と曲げモーメント $M(x)$ を求め,
        せん断力図（SFD）と曲げモーメント図（BMD）を示せ.

  \item たわみ $v(x)$ に関するはりの基礎式を示し,
        たわみ角とたわみ曲線を求め, C 点におけるたわみを求めよ.

  \item C 点におけるたわみをカスチリアーノの定理（エネルギー法）を用いて求めよ.
\end{enumerate}
